% Separate post-CP28 candidate; the final abelian-image branch remains unresolved.
\subsection{The next wild image and its actual integral Cartan}

Retain the hypotheses of B79: \(p\ge5\), \(q=p^2\), \(N=q-1\),
\(e_K=p(p-1)\), algebraically closed residue field and \(\mu_p\subset K\).
The common residual inertia representation is irreducible, and the
comparison, when used, is an actual \(p^3\)-torsion isomorphism.
Let \(L/K\) have tame degree \(N\), let \(F/L\) be the common full
third-level division field, and suppose \(|W|=|\operatorname{Gal}(F/L)|=p^4\).
Use the actual matrix spaces \(W_1,W_2\) of (B-RD79.2).
For each formal model, the primitive, second-level and first-level
coordinate valuations in \(F\) are respectively
\[
 A=p(p-1),\qquad qA,\qquad q^2A,
                 \qquad v_F(p)=Nq^2A.                 \tag{B-RD80.1}
\]
A tame complement of order \(N\) exists in the actual finite
Galois group: a lift of a tame generator has order \(Np^a\),
and taking its \(p^a\)th power gives an order-\(N\) lift of
another tame generator. Write \(T\) for its action on the residual module,
and identify its centralizer with \(\mathbf F_q\). Denote the
Cartan and anti-Cartan planes in its matrix algebra by
\[
 Z=\{z\mapsto az:a\in\mathbf F_q\},\qquad
 U=\{z\mapsto az^p:a\in\mathbf F_q\}.
\]
Thus tame conjugation fixes \(Z\) and acts on \(U\) through
multiplication by a root of order \(h\mid p+1\), where \(h>1\)
is the residual projective order.

\begin{lemma}[A first-coefficient comparison at either displacement level]
\label{lem:B-RD80-first}
\status{Candidate; complete proof supplied, independent audits pending.}
Fix a Honda coordinate in one model and fixed residue representatives
in a coefficient ring of \(K\). For each \(z\in\mathbf F_q^\times\),
choose one primitive third-level point with leading coefficient \(dz\),
and keep these choices fixed. Let \(g\in W\) have lower break \(b\),
and suppose its displacement on all these points has valuation
\(D\), where \(D=qA\) or \(D=q^2A\). Suppose its leading displacement
map is a nonzero scalar multiple of either \(z\) or \(z^p\).
Put \(b_0=(D-A)/p\).
If \(b<b_0\), the first nonzero coefficient of degree prime to \(p\)
in the fixed point expansions below \(D\) has degree
\[
             j=A+(p-1)b,
 \qquad f_j(z)=Cz\,d_g^{p-1},\qquad C\ne0,           \tag{B-RD80.2}
\]
where \(d_g\) is the leading uniformizer displacement, and the
constant \(C\) depends on the fixed break and model, not on \(g\).
In particular \(p\nmid b\).
If \(b>b_0\), that first degree is \(j=D-b\), and its coefficient
is proportional to the leading displacement map divided by \(d_g\).
These conclusions concern the same fixed point expansions for every
chosen group element; no independence from other lifts modulo \(p\)
is asserted.
\end{lemma}
\begin{proof}
Work through degree \(D\); the coefficient ring is of characteristic
\(p\) in this truncated calculation because \(v_F(p)>D\).
A monomial of degree \(j\) first changes at degree
\(j+p^{v_p(j)}b\). The initial degree is \(A\), with \(v_p(A)=1\).
Among all degrees divisible by \(p\), it is the unique one whose
first change can occur at \(A+pb\); every later such degree changes
strictly later. If \(b<b_0\), this initial change must cancel.
No earlier prime-to-\(p\) degree can occur, since it would change
alone even earlier; none later can cancel. The required first degree
is therefore \(j=A+(p-1)b\), prime to \(p\). The leading cancellation is
\[
       (A/p)dz\,d_g^p+jf_j(z)d_g=0.
\]
Dividing by \(jd_g\) proves (B-RD80.2), including \(p\nmid b\).
If \(b>b_0\), every degree divisible by \(p\) changes after \(D\).
The first prime degree must instead be \(D-b\), and its linear
binomial change gives the stated coefficient. Any impossible
index or zero coefficient would contradict the prescribed nonzero
displacement. These are the full fixed-coefficient calculations of
B56--B57 with a displayed arbitrary target \(D\); they also apply
to a group element of order \(p^2\).
\end{proof}

\begin{proposition}[The three-dimensional terminal layer is excluded]
\label{prop:B-RD80-three-terminal}
\status{Candidate; complete proof supplied, independent audits pending.}
If \(\dim W_1=1\) and \(\dim W_2=3\), two actual compared models
cannot have different finite special characters.
\end{proposition}
\begin{proof}
For a nonzero matrix in either \(Z\) or \(U\), every primitive
point is moved by a nonzero first-level point. Apply the lemma
with \(D=q^2A\) to the corresponding terminal elements.
Here \(b_0=(q^2-1)(p-1)\) is a multiple of \(N\).
Choose an exact tame eigen-uniformizer of \(F\), by averaging
an arbitrary uniformizer. Its eigenvalue \(\eta\) has order \(N\).
On a nonzero ramification-graded piece of lower break \(b\),
tame conjugation acts by \(\eta^b\) (reversing both conventions
has no effect on the following tests). The leading coefficient
map on a subgroup with one common break is injective and additive.

If \(h>2\), the anti-Cartan plane is irreducible over \(\mathbf F_p\).
Any three-dimensional tame-invariant plane consequently contains
all of \(U\), with one common break \(b\); moreover \(\eta^b\ne1\),
so \(b\ne b_0\). If \(b<b_0\), fix any \(z\ne0\) in (B-RD80.2)
and vary the nonzero elements of \(U\). Its two-dimensional
leading-coefficient space would have all its nonzero elements
satisfy the same equation \(X^{p-1}=c\), which has at most
\(p-1\) roots. This contradicts its \(p^2-1\) nonzero elements.
If \(b>b_0\), the high-break formula and one fixed element show
that the displacement parameters \(\beta_i(g)\) of each model
are fixed nonzero multiples of the same \(d_g\). The parameters
run through \(\mathbf F_q\). The Frobenius change between the
two leading orientations, proved at this absolute index in B79,
requires \(\beta_2(g)=a^{1-p}\beta_1(g)^p\) for the fixed
nonzero scalar \(a\) of that intertwiner, again impossible on the
whole field.
This proves the claim for \(h>2\).

Suppose \(h=2\). Tame conjugation is \(-1\) on \(U\), so every
anti-Cartan break satisfies \(b\equiv N/2\pmod N\), whereas
Cartan breaks satisfy \(b\equiv0\pmod N\). If all of \(U\)
is present but has two breaks \(b_1<b_2\), then \(b_1\ne b_0\).
For \(b_1<b_0\), the first prime degree is
\(j=A+(p-1)b_1\), with coefficient a nonzero multiple of \(z\).
Under an element at \(b_2\), its change at \(j+b_2\) is uniquely
first: it precedes \(A+pb_2\) and every later prime degree.
This would give a Cartan-linear displacement, contrary to an
anti-Cartan matrix. For \(b_1>b_0\), the first prime degree is
\(D-b_1\), so all its changes at \(b_2\), and all divisible
ones, occur after \(D\), again impossible. Thus a full \(U\)
has one break, and the previous single-break argument excludes it.

The remaining three-dimensional shape is \(Z\oplus U_1\), with
\(U_1\) an anti-Cartan line. Let \(b\) be its nonzero break.
It is neither a Cartan break nor \(b_0\). If \(b>b_0\), the first
prime coefficient is a multiple of \(z^p\), at \(j=D-b\).
A Cartan break \(c<b_0\) would require that same first coefficient
to be a multiple of \(z\). For \(b_0\le c<b\), the degree-\(j\)
term changes strictly before \(D\), while the divisible terms
do not; for \(c>b\), every change is later than \(D\).
The equality \(c=b\) is forbidden by tame characters. Hence
\(b<b_0\). Its first prime degree is now \(j=A+(p-1)b\), with
coefficient proportional to \(z\). A Cartan break smaller than
\(b\) would force an earlier first prime degree; equality is
again impossible. Each Cartan element must therefore have
\[
                      c=D-j=pb_0-(p-1)b>b_0.
\]
All nonzero elements of \(Z\) have this same break. The high-break
formula then makes their multiplication parameters fixed multiples
of the common \(d_g\) in both models. Frobenius orientation on
this full field again contradicts scalar proportionality. This
exhausts every possible three-dimensional terminal plane.
\end{proof}

\begin{lemma}[A correct integral-lift test for the two-dimensional layers]
\label{lem:B-RD80-plane}
\status{Candidate; complete proof supplied, independent audits pending.}
If \(\dim W_1=\dim W_2=2\), their common matrix plane is a
nonsplit Cartan over \(\mathbf F_p\). The full wild group is
abelian, isomorphic to \((\mathbf Z/p^2)^2\).
\end{lemma}
\begin{proof}
The inclusion from (B-RD79.3) gives \(W_1=W_2=V\).
The exact commutator formula modulo \(p^3\) shows that
\([V,V]\subset V\). Thus \(V\) is a two-dimensional matrix
Lie algebra normalized by the irreducible residual tame matrix.
If it is nonabelian, its characteristic derived line has a generator
\(E\ne0\) and one may choose \(H\) with \([H,E]=E\).
Then \(\operatorname{tr}E=0\), and
\([H,E^2]=2E^2\) gives \(\operatorname{tr}E^2=0\).
Since \(p\ne2\), \(\det E=-\operatorname{tr}(E^2)/2=0\).
So \(E\) is a nonzero nilpotent and its kernel is an
\(\mathbf F_p\)-line preserved by the tame normalizer,
contradicting irreducibility.

The plane is therefore abelian. Choose a nonscalar element.
Its two-dimensional centralizer equals \(V\), and is either
a nonsplit Cartan, a split Cartan, or the dual-number algebra.
In the last case the unique nilpotent line and its kernel are
again preserved by the tame normalizer, a contradiction.

For the split Cartan use an eigenbasis so that \(V\) is the
entire diagonal plane. Choose an actual element
\(g=1+pE_{22}+p^2B\in W\). Multiplication by an actual terminal
diagonal matrix permits \(B_{11}=0\) modulo \(p\).
Choose a primitive vector
\[
                 v=e_1+p w,\qquad w_2=-B_{21}\pmod p.
                                                     \tag{B-RD80.3}
\]
Then \((E_{22}+pB)v=0\pmod{p^2}\), so \(gv=v\pmod{p^3}\).
But \(g^p=1+p^2E_{22}\ne1\), so \(g\) has order \(p^2\).
Its stabilizer on this actual primitive point has order at least
\(p^2\), giving an orbit of size at most \(p^2\). This contradicts
its degree-\(p^3\) valuation denominator from B79.
The choice of the integral lift in (B-RD80.3) is essential: no
bound by a sum of residual evaluation ranks is used.

Only a nonsplit Cartan remains. Its matrices commute, so the
actual commutator formula gives \([W,W]=1\) modulo \(p^3\).
Taking \(p\)th powers maps onto the terminal plane of order
\(p^2\). An abelian group of order \(p^4\), exponent at most
\(p^2\), with a \(p\)th-power subgroup of order \(p^2\), is
\((\mathbf Z/p^2)^2\), as asserted.
\end{proof}

\begin{lemma}[The terminal Cartan is the original residual Cartan]
\label{lem:B-RD80-orientation}
\status{Candidate; complete proof supplied, independent audits pending.}
In the situation of the preceding lemma, \(V\) equals the
centralizer of the original irreducible residual tame matrix.
\end{lemma}
\begin{proof}
Write \(E'\) for the nonsplit Cartan \(V\). The normalizer of
\(E'\) is \(E'^\times\rtimes\operatorname{Gal}(E'/\mathbf F_p)\):
conjugation restricts to a field automorphism, and its kernel is
field multiplication. If the tame matrix lies in \(E'^\times\),
its irreducibility identifies its centralizer with \(E'\), as desired.
Suppose instead it lies in the other coset. Its square is scalar,
and its conjugation on \(E'\) has eigenspaces \(\mathbf F_p I\)
and a trace-zero line. Relative to the original residual field,
these are a Cartan scalar line and an anti-Cartan line.

The actual tame action on \(W\simeq(\mathbf Z/p^2)^2\) has order
prime to \(p\). Its square is the identity modulo \(pW\).
The reduction kernel in its automorphism group is a \(p\)-group,
so its square is the identity on \(W\) itself. Since two is a
unit, the exact \(\pm1\) projectors give two rank-one summands.
Choose generators \(g_+\) and \(g_-\), each of order \(p^2\),
with tame conjugation respectively \(g_+\mapsto g_+\) and
\(g_-\mapsto g_-^{-1}\). Their first matrices are respectively
nonzero scalar and anti-Cartan matrices, both invertible.
Their displacements on every primitive point have valuation \(qA\).

Apply the first-coefficient lemma with \(D=qA\) and
\(a_0=(q-1)(p-1)\). Write \(a\) for the lower break of \(g_-\)
and \(c\) for that of \(g_+\). Tame characters give
\(a\equiv N/2\pmod N\), \(c\equiv0\pmod N\).
Thus \(a\ne c\) and \(a\ne a_0\).
If \(a>a_0\), the first prime coefficient has degree \(qA-a\)
and is a nonzero multiple of \(z^p\). A break \(c<a_0\)
would force that first coefficient to be a multiple of \(z\);
for \(a_0\le c<a\) it would change strictly before \(qA\),
with no divisible term available to cancel; for \(c>a\) all
changes would occur later than \(qA\). Thus \(a<a_0\).
The first prime degree is
\(j=A+(p-1)a\), and its coefficient is a nonzero multiple of \(z\).

Now use the actual element \(g_-^p\), whose nonzero terminal
matrix is the same anti-Cartan matrix. Its lower break \(b\)
is strictly greater than \(a\): the first ramification coefficient
of a product is additive, so the first coefficient of the \(p\)th
power vanishes in the residue field. Its displacement must have
valuation \(q^2A\) with leading coefficient proportional to \(z^p\).
But the already fixed degree-\(j\) coefficient changes at
\(j+b<A+pb\), before every divisible term and every later
prime term. It is uniquely first and is proportional to \(z\).
If its degree differs from \(q^2A\), valuations already contradict;
if it equals \(q^2A\), the independence of \(z\) and \(z^p\) as
functions on \(\mathbf F_q\) contradicts the required leading map.
This excludes the noncentral normalizer case and proves the lemma.
\end{proof}

\begin{theorem}[The fourth-order wild branch has an actual abelian image]
\label{thm:B-RD80}
\status{Candidate; complete proof supplied, independent audits pending.}
If the two finite special characters differ and their common
third-level wild image has order \(p^4\), then
\[
 W_1=W_2=Z,
 \qquad W\simeq(\mathbf Z/p^2)^2,
                                                     \tag{B-RD80.4}
\]
and the full common third-level inertia image is abelian.
More precisely its Tate module modulo \(p^3\) is free of rank one
over the unramified quadratic algebra \(\mathcal O_E/p^3\),
its full image acts by multiplication, and its wild image is exactly
\[
                 (1+p\mathcal O_E)/(1+p^3\mathcal O_E).
                                                     \tag{B-RD80.5}
\]
This is a statement about the finite-level action, not a CM
endomorphism assertion for the full BT group.
\end{theorem}
\begin{proof}
As in B79, \(W_1\ne0\), \(W_1\subset W_2\), and their dimensions
sum to four. The only possibilities are \((1,3)\) and \((2,2)\).
The former is excluded for a different-character pair by the
three-dimensional terminal proposition. The next two lemmas give
(B-RD80.4) in the latter case.

The tame complement acts trivially on \(W/pW=W_1=Z\).
Its finite prime-to-\(p\) action on \(W\simeq(\mathbf Z/p^2)^2\)
therefore is trivial, again because the reduction kernel in its
automorphism group is a \(p\)-group. Thus the complement centralizes
\(W\), and the full finite inertia image is abelian.

Let \(\widetilde T\) be the actual third-level tame matrix.
Its residual quadratic polynomial is irreducible and separable.
The algebra it generates over \(\mathbf Z/p^3\) is therefore the
unramified quadratic algebra \(\mathcal O_E/p^3\), and the free
rank-two underlying module is free of rank one over that algebra.
Every matrix commuting with \(\widetilde T\) is multiplication
by an element of it, as follows either from the basis
\(v,\widetilde T v\) or the centralizer equations.
The entire image commutes with \(\widetilde T\), so it lies in
these units. Its wild subgroup lies in their principal units,
whose order is \(p^4\). Equality of orders proves (B-RD80.5).
\end{proof}

\begin{remark}[What the finite Cartan does not prove]
\label{rem:B-RD80-scope}
\status{Candidate scope; the abelian fourth-order branch remains open.}
This narrows the \(p^4\) branch left by B79 to a full finite
Cartan image. It neither excludes that branch nor identifies a
finite-level quadratic action with finite-extension endomorphisms
of the whole group. B78's CM theorem therefore cannot be applied
without an additional argument. Wild orders \(p^5,\ldots,p^8\)
and reducible residual representations remain open as well.

The earlier exploratory orbit-rank shortcut was withdrawn and its
exact note retained. The split-Cartan argument here instead constructs
an actual primitive lift with an order-\(p^2\) stabilizer. No assertion
that all fibre motions come from \(W_2v\) is used.
\end{remark}
